\begin{lemma}[A maximizer on the latitude set]
Let $n\ge2$, let $e\in\mathbb{R}^n$ with $\lVert e\rVert=1$, where
$\mathbb{R}^n$ carries the Euclidean inner product, and let
$f:\mathbb{R}^n\to\mathbb{R}$ be continuous. Then there is
$v_0\in\mathbb{R}^n$ with
\[
  \lVert v_0\rVert=1,
  \qquad
  \langle e,v_0\rangle=0,
  \qquad
  f(v)\le f(v_0)\ \text{ for every }v\text{ with }\lVert v\rVert=1,\ \langle e,v\rangle=0 .
\]
\end{lemma}
\begin{proof}
Write $S=\{v\in\mathbb{R}^n:\lVert v\rVert=1\ \text{and}\ \langle e,v\rangle=0\}$
for the latitude set.

\emph{Step 1: $S$ is compact.} We have
$S=\bigl\{v:\lVert v\rVert=1\bigr\}\cap\bigl\{v:\langle e,v\rangle=0\bigr\}$ by
definition. The first factor, the unit sphere, is compact by the background
proposition on the unit sphere. The second factor is closed: the map
$v\mapsto\langle e,v\rangle$ is continuous (it is linear on a
finite-dimensional space, being the inner product with a fixed vector) and the
set is the preimage of the closed set $\{0\}$ under it. The intersection of a
compact set with a closed set is compact, so $S$ is compact.

\emph{Step 2: $S$ is nonempty.} Since $\lVert e\rVert=1$ we have $e\ne0$, so the
line $\mathbb{R}e$ spanned by $e$ has dimension $1$. For a subspace $K$ of a
finite-dimensional inner product space one has
$\dim K+\dim K^{\perp}=\dim\mathbb{R}^n=n$; applying this with $K=\mathbb{R}e$
gives $\dim(\mathbb{R}e)^{\perp}=n-1\ge1$, using $n\ge2$. A subspace of positive
dimension is not the zero subspace, so there is $w\in(\mathbb{R}e)^{\perp}$ with
$w\ne0$. By definition of the orthogonal complement, $w$ is orthogonal to every
element of $\mathbb{R}e$, in particular $\langle e,w\rangle=0$. Since $w\ne0$ we
have $\lVert w\rVert\ne0$, so $v=\lVert w\rVert^{-1}w$ is defined and satisfies
\[
  \lVert v\rVert=\lVert w\rVert^{-1}\lVert w\rVert=1,
  \qquad
  \langle e,v\rangle=\lVert w\rVert^{-1}\langle e,w\rangle=0 ,
\]
using that the norm is absolutely homogeneous and the inner product is linear in
its second argument. Hence $v\in S$ and $S\ne\emptyset$.

\emph{Step 3: conclusion.} By Steps 1 and 2, $S$ is a nonempty compact set, and
$f$ is continuous, hence continuous on $S$. A continuous real-valued function on
a nonempty compact set attains a maximum there, so there is $v_0\in S$ with
$f(v)\le f(v_0)$ for every $v\in S$. Unfolding the definition of $S$, this
$v_0$ satisfies $\lVert v_0\rVert=1$, $\langle e,v_0\rangle=0$, and
$f(v)\le f(v_0)$ for every $v$ with $\lVert v\rVert=1$ and
$\langle e,v\rangle=0$, which is the assertion.
\end{proof}
